\paragraph{What this changes about external validation.}
The standard package of evidence --- an AUROC obtained in a second country ---
is not wrong, it is answering a different question from the one deployment
poses.  Our transfer matrix reproduces the reassuring AUROC and shows
calibration failing behind it.  Any use of a model that is not pure ranking
depends on the half that goes unreported.  We propose that a deployment dossier
state, per site and per label, the calibration ratio, the decomposition, and the
label count required to close the gap --- all of which are computable with the
released toolbox, and two of which need no target labels at all.

\paragraph{What a hospital should actually do.}
The workflow the results support has three steps and a decision point.  Run the
unlabeled goodness-of-fit test on your own recordings.  If it does not reject,
the prevalence correction is licensed and free.  If it does, the shift at your
site is not reducible to prevalence: use the minimax correction, which damps
itself by how little the data pins the prevalence, and budget for labels using
the plug-in formula, whose only unknown is your model's mean predictive
variance on your own unlabeled data.  The number of labels this implies is in
the hundreds, not the tens of thousands a retraining effort would need.

\paragraph{Implications for LMIC deployment.}
The asymmetry here is favourable and worth stating plainly.  The component of
the cliff that is free to fix is the prevalence component, and prevalence is
exactly what differs most between a tertiary centre in a high-income country
and a primary-care setting elsewhere.  A site with little capacity for chart
review is disproportionately likely to be in the regime where it needs none.
The corollary is equally plain: that site still needs the unlabeled test, since
without it the free correction is a coin flip that can make matters worse.

\paragraph{Limitations.}
\emph{Label provenance.}  Most of these cohorts take their labels from an
interpretation of the ECG, so a cross-country difference in labels can always
be partly a difference in reading practice rather than in patients.  The
EchoNext arm, whose labels come from echocardiography, is the control for this,
and it is the one arm we regard as decisive on the point.

\emph{The Korean cohort.}  Its inclusion requires institutional approval and a
data transfer agreement.  The federated protocol removes the technical
obstacle, not the governance one.

\emph{$\gamma$ is an assumption.}  Both the partial-identification interval and
the conformal inflation depend on a concept-shift budget that cannot be
estimated from the data it is used on.  We report sweeps rather than single
numbers, and $\gamma_{\min}$ as the data's own floor; a reader who disbelieves
our $\gamma$ can read their own off the sweep.

\emph{Conditional versus unconditional test size.}  The unlabeled test is
calibrated by parametric bootstrap and has close to nominal unconditional size.
In deployment, though, every site uses \emph{one} source operator estimated by
whoever shipped the model, and conditional on a particular such operator the
size can be materially higher.  This is why we recommend the minimax
correction, which involves no binary decision, over the gate.  A vendor
shipping a model for this workflow should estimate the source operator on as
much data as possible and report its conditioning.

\emph{Simulation versus cohorts.}  The theorems are validated on a generator
with known ground truth because the split they concern is unobservable in real
data.  The empirical claims about transfer are separate and belong to the
cohorts; we have not used the simulator to support them, and the two are never
pooled in one table.

\paragraph{What would falsify this.}
A site with a large prevalence shift, a non-rejecting goodness-of-fit test, and
a free correction that fails to recover most of its calibration error would
contradict Theorem~\ref{thm:decomp}(c) as applied here.  A site where the
measured label requirement materially exceeds the plug-in bound at its measured
$\Gamma$ would contradict Theorem~\ref{thm:samplesize}(a).  Both are checkable
with the released code on any cohort a reader has access to, which is the point
of releasing it.
